

\zadanieM{1807}{Dane są dwa ciągi arytmetyczne, $a_1,$ $a_2,$ $\ldots$ i~$b_1,$ $b_2,$ $\ldots$, składające się z~liczb całkowitych dodatnich. Wiadomo, że $a_1 = b_1$ oraz dla każdej liczby całkowitej $n\geq 1$ liczby $a_n$ i~$b_n$ dają równe reszty z~dzielenia przez $n.$ Udowodnić, że te ciągi są identyczne.}{Skoro pierwsze wyrazy danych ciągów
są takie same, wystarczy udowodnić równość ich różnic $a$ i~$b$. Z~warunków
zadania wynika, że dla dowolnej liczby całkowitej dodatniej $n$
liczba
\begin{align*}
  b_{n}-a_{n}& =(b_{1}+(n-1)b)-(a_{1}+(n-1)a)\\
             &=(n-1)(b-a)
\end{align*}
jest podzielna
przez $n,$ co oznacza, że liczba $b-a$ jest podzielna przez $n.$
Stwierdziliśmy więc, że liczba całkowita $b-a$ jest podzielna przez dowolną
liczbę całkowitą dodatnią, a~to jest możliwe tylko wtedy, gdy $b-a=0.$}
\zadanieM{1808}{Dane są dwa ciągi liczb całkowitych dodatnich $x_{1},$ $x_{2},$ $\ldots,$ $x_{m}$ oraz $y_{1},$ $y_{2},$ $\ldots,$ $y_{n},$ których sumy $x_{1}+x_{2}+\ldots +x_{m}$ i~$y_{1}+y_{2}+\ldots +y_{n}$ są równe $s<nm.$ Udowodnić, że w~równości 
$$x_{1}+x_{2}+\ldots +x_{m}=y_{1}+y_{2}+\ldots +y_{n}$$ można wykreślić część składników, zachowując jej prawdziwość.}{Rozważmy okrąg o~długości $s$ i~podzielmy go na równe łuki o~długości $1.$ Zaznaczmy na czerwono $m$ punktów podziału okręgu na łuki o~długościach $x_{1},$ $x_{2},$ $\ldots,$ $x_{m}$ i~kolorem niebieskim $n$ punktów podziału okręgu na łuki o~długościach $y_{1},$ $y_{2},$ $\ldots,$ $y_{m}.$ Na podstawie zadania \href{https://www.deltami.edu.pl/2025/01/zadania/}{\textbf{M 1806}} z~poprzedniego numeru pewne dwa łuki, $A_{i_1} A_{i_2}$ i~$B_{j_{1}} B_{j_2},$ są równej długości. Usuwając te długości z~równości, zachowujemy jej prawdziwość. }
\zadanieM{1809}{\marg[-40]{\begin{tikzpicture}[scale=1.3]
%%% definicje
\tkzDefPoints{0/0/A, 3/0/B, 1.3/.5/C, 1.7/3/D}
\tkzDefBarycentricPoint(C=1.4,D=1)\tkzGetPoint{X}
\tkzDefBarycentricPoint(X=1.5,B=1)\tkzGetPoint{A'}
\tkzDefBarycentricPoint(X=1.3,A=1)\tkzGetPoint{B'}
%%% rysunki
\tkzDrawSegments[thick](A,B B,C C,A A,D B,D C,D A,A' B,B')
\tkzDrawSegments[dashed](A',B A',C A',D B',A B',C B',D)
\tkzMarkAngles[size=2mm,mark=|,mksize=2pt](A,B',C C,A',B)
\tkzMarkAngles[size=1.5mm,mark=|,mksize=2pt](D,B',A B,A',D)
%%% etykiety
\tkzLabelPoints[below,font=\scriptsize](A,B,C)
\tkzLabelPoints[above,font=\scriptsize](D)
\node[xshift=-1.8mm,yshift=1.5mm] at (A') {\scriptsize $A'$};
\node[xshift=2mm,yshift=1.5mm] at (B') {\scriptsize $B'$};
\end{tikzpicture}}W~czworościanie $ABCD$ na ścianach $BCD$ i~$ACD$ znajdują się, odpowiednio, punkty $A'$ oraz $B'$ takie, że $$\angle AB'C = \angle AB'D = \angle BA'C = \angle BA'D = 120^{\circ}.$$ Wiadomo, że proste $AA'$ i~$BB'$ przecinają się. Udowodnić, że odległości punktów $A'$ oraz $B'$ od prostej $CD$ są równe.}{Z~warunków zadania wynika, że punkty $A,$ $A'$, $B,$ $B'$ leżą w~tej samej
płaszczyźnie, zatem proste $BA'$ i~$AB'$ przecinają krawędź $CD$ w~jednym
punkcie $X.$ Z~równości kątów wynika, że proste te są dwusiecznymi kątów,
odpowiednio, $CA'D,$ $CB'D.$ Zatem z~twierdzenia o~dwusiecznej dostajemy
\[\frac{CA'}{A'D}= \frac{CX}{XD} = \frac{CB'}{B'D},\]
a~ponieważ $\angle CA'D =
\angle CB'D = 120^{\circ},$ trójkąty $CA'D$ i~$CB'D$ są podobne. Ponieważ $CD$
jest wspólnym bokiem tych trójkątów, są one przystające. Oczywiście w~tych przystających trójkątach odpowiadające im wysokości z~wierzchołków $A'$ i~$B'$ są równe.\medskip

\centerline{\begin{tikzpicture}[scale=1.3]
%%% definicje
\tkzDefPoints{0/0/A, 3/0/B, 1.3/.5/C, 1.7/3/D}
\tkzDefBarycentricPoint(C=1.4,D=1)\tkzGetPoint{X}
\tkzDefBarycentricPoint(X=1.5,B=1)\tkzGetPoint{A'}
\tkzDefBarycentricPoint(X=1.3,A=1)\tkzGetPoint{B'}
%%% rysunki
\tkzDrawSegments[thick](A,B B,C C,A A,D B,D C,D)
\tkzDrawSegments(A,A' B,B')
\tkzDrawSegments[magenta](A,X B,X)
\tkzDrawSegments[thick,magenta](B',C B',D A',C A',D)
%%% etykiety
\tkzLabelPoints[font=\scriptsize](A,B,C)
\tkzLabelPoints[above,font=\scriptsize](D)
\node[xshift=0,yshift=-6pt] at (A') {\scriptsize $A'$};
\node[xshift=0,yshift=-6pt] at (B') {\scriptsize $B'$};
\node[xshift=5pt,yshift=2pt] at (X) {\scriptsize $X$};
\end{tikzpicture}}}\goodbreak

\zadanieF{1113}{Planeta o~masie $m$ obiega S\l{}o\'{n}ce po orbicie eliptycznej.
W~aphelium jej odleg\l{}o\'{s}\'{c} od S\l{}o\'{n}ca wynosi $A$, a~w~perihelium $a$.
Ile wynosi moment p\k{e}du, $L$, tej planety? Masa S\l{}o\'{n}ca r\'{o}wna jest $M$, a~sta\l{}a~grawitacji $G$.}{Skorzystamy z~praw: zachowania energii i~momentu
p\k{e}du. W~perihelium i~w~aphelium pr\k{e}dko\'{s}\'{c} planety jest prostopad\l{}a~do odcinka \l{}\k{a}cz\k{a}cego \'{s}rodek planety ze \'{s}rodkiem S\l{}o\'{n}ca. Oznaczmy pr\k{e}dko\'{s}ci w~aphelium i~w~perihelium, odpowiednio, jako $v_A$ i~$v_a$. Mamy:
\begin{displaymath}
 L = mAv_A = mav_a,
\end{displaymath}
czyli: $v_A = L/(mA)$ i~$v_a = L/(ma)$,
oraz
\begin{displaymath}
 -\frac{GMm}{A} + \frac{mv_A^2}{2} = -\frac{GMm}{a} + \frac{mv_a^2}{2}.
\end{displaymath}
Po skorzystaniu ze zwi\k{a}zk\'{o}w~pr\k{e}dko\'{s}ci z~momentem p\k{e}du otrzymujemy r\'{o}wnanie:
\begin{displaymath}
 -\frac{GMm}{A} + \frac{L^2}{2mA^2} = -\frac{GMm}{a} + \frac{L^2}{2ma^2}.
\end{displaymath}
Po prostych przekszta\l{}ceniach otrzymujemy odpowied\'{z}:
\begin{displaymath}
 L = m\sqrt{\frac{2GMAa}{A+a}}.
\end{displaymath}}
\zadanieF{1114}{W~opisanym dalej uk\l{}adzie stoj\k{a}ca fala d\'{z}wi\k{e}kowa mo\.{z}e spe\l{}nia\'{c} rol\k{e} siatki dyfrakcyjnej dla \'{s}wiat\l{}a. Kuweta
o~przezroczystych, p\l{}askich, r\'{o}wnoleg\l{}ych \'{s}ciankach jest wype\l{}niona ciecz\k{a}.  R\'{o}wnolegle do \'{s}cianek kuwety w~cieczy wzbudzona jest stoj\k{a}ca fala ultrad\'{z}wi\k{e}kowa o~d\l{}ugo\'{s}ci $\Lambda$. Wi\k{a}zka \'{s}wiat\l{}a~laserowego o~d\l{}ugo\'{s}ci fali $\lambda$ pada prostopadle do \'{s}cianek kuwety i~do wektora falowego ultrad\'{z}wi\k{e}ku. Jaki warunek spe\l{}niaj\k{a} powstaj\k{a}ce za kuwetą pr\k{a}\.{z}ki interferencyjne \'{s}wiat\l{}a?}{Zmianom ci\'{s}nienia towarzysz\k{a} zmiany pr\k{e}dko\'{s}ci \'{s}wiat\l{}a~w~cieczy. W~zwi\k{a}zku z~tym kolejne zg\k{e}szczenia i~rozrzedzenia cieczy wywo\l{}ane fal\k{a} akustyczn\k{a} b\k{e}d\k{a} dzia\l{}a\l{}y jak szczeliny siatki dyfrakcyjnej. Przyjmijmy, \.{z}e stoj\k{a}ca fala ultrad\'{z}wi\k{e}kowa wzbudzana jest w~kierunku $x$. Przesuni\k{e}cia cz\k{a}stek cieczy w~kierunku $x$ (d\'{z}wi\k{e}k to fala pod\l{}u\.{z}na) opisywane s\k{a} w\'{o}wczas wzorem (suma dw\'{o}ch fal biegn\k{a}cych o~przeciwnych zwrotach wektorów falowych):
\begin{align*}
  u(x,t) &= u_0 (\sin(\omega t - kx) - \sin(\omega t + kx))\\
         &= 2u_0\cos(\omega t)\sin(kx).
\end{align*}
W~powy\.{z}szych wzorach $k = 2\pi/\Lambda$ oznacza d\l{}ugo\'{s}\'{c} wektora falowego,
a $\omega$ ko\l{}ow\k{a} cz\k{e}sto\'{s}\'{c} fali. Zbadajmy rozk\l{}ad pr\k{e}dko\'{s}ci:
\begin{displaymath}
 v(x,t) = -2u_0\omega \sin(\omega t)\sin(kx).
\end{displaymath}
W\k{e}z\l{}y fali ($v(x,t) = 0$) pojawiaj\k{a} si\k{e} w~punktach, dla kt\'{o}rych
$kx = n\pi$ dla ca\l{}kowitych $n$. Nie wszystkie te w\k{e}z\l{}y s\k{a} r\'{o}wnowa\.{z}ne, bo je\'{s}li przyjmiemy, \.{z}e $-2u_0\omega \sin(\omega t) > 0$, to
w~w\k{e}z\l{}ach odpowiadaj\k{a}cych nieparzystym $n$ cz\k{a}stki z~obu stron poruszaj\k{a} si\k{e} w~kierunku w\k{e}z\l{}a~(zmiana znaku~$v$ z~dodatniego na ujemny),
a~w~w\k{e}z\l{}ach o~parzystym~$n$ oddalaj\k{a} si\k{e}. Odleg\l{}o\'{s}ć kolejnych rozrzedze\'{n} cieczy (,,szczelin'' siatki) wynosi wi\k{e}c $\Lambda$ i~jest taka sama jak
w~przypadku rozpraszania na fali biegn\k{a}cej. Kolejne pr\k{a}\.{z}ki interferencyjne
pojawiaj\k{a} si\k{e} wi\k{e}c pod k\k{a}tami $\alpha$ do wi\k{a}zki padaj\k{a}cej  spe\l{}niaj\k{a}cymi warunek:
\begin{displaymath}
 \Lambda \sin(\alpha) = m\lambda
\end{displaymath}
z~całkowitymi warto\'{s}ciami $m$. 


Uwaga: Cz\k{e}sto\'{s}ci ultrad\'{z}wi\k{e}k\'{o}w, rz\k{e}du MHz, s\k{a} tak ma\l{}e
w~por\'{o}wnaniu z~cz\k{e}sto\'{s}ciami fal \'{s}wietlnych, \.{z}e \'{s}wiat\l{}o~rozprasza si\k{e} na fali ultrad\'{z}wi\k{e}kowej jak na nieruchomej siatce dyfrakcyjnej.}
